% problem-000351 5 超临界水热力学
\section*{5. 超临界水热力学}
\textit{下列为分问。}

\subsection*{5-1}
画出⽔的压强-温度(�-�)相图（⽰意图），标注⽓相、液相、固相和超临界⽔所在的区域。

\subsection*{5-2}
计算液态⽔在 $9 0 ^ { \circ } \mathrm { C }$ 和 0.1 MPa 时的 $) K _ { \mathrm { w } } ^ { \ominus } ( \Delta _ { \mathrm { r } } H _ { \mathrm { m } } ^ { \ominus }$ 和 $| \Delta _ { \mathrm { { r } } } S _ { \mathrm { { m } } } ^ { \ominus }$ 可视为常数，相关数据⻅表1）。

\subsection*{5-3}
有⼈研究了⼄酸⼄酯在23\~30 MPa、 $2 5 0 { \sim } 4 0 0 \ ^ { \circ } \mathrm { C }$ 和没有任何其他外加物的条件下的⽔解动⼒学，并提出两种可能的机理。

机理 1：

$
\begin{array}{r l} & \mathrm {CH_ {3} COOH\xrightarrow [ k _ {- 1} ]{k_ {1}} CH_ {3} COO^ {-} + H^ {+} \qquad K_ {a}} \\ & \mathrm {H_ {3} C-OC(OC_ {2} H_ {5}) + H^ {+} \xrightarrow [ k _ {- 2} ]{k_ {2}} \left[ H_ {3} C-OC(OC_ {2} H_ {5}) \rightleftharpoons H_ {3} C-OC(OC_ {2} H_ {5})} \quad K _ {2} \\ & \mathrm {H_ {3} C-OC(OC_ {2} H_ {5}) + 2 H_ {2} O\xrightarrow [ k _ {3} ]{k _ {3}} \left[ H_ {3} C-OC(OC_ {2} H_ {5}) \rightleftharpoons H_ {3} O^-H]} \\ & \left[ \begin{array}{c} H \\ H _ {3} C - \mathrm{O} - H \\ \mathrm{H} _ {3} C - \mathrm{O} - H \\ \mathrm{OC} _ {2} \mathrm{H} _ {5} \end{array} \right] \xrightarrow [ \text {fast} ]{k _ {4}} H _ {3} C - \mathrm{O} - H + H _ {3} O ^ {+} \\ & \mathrm {H_ {3} C-OC(OC_ {2} H_ {5})\xrightarrow [ \text {fast} ]{k_ {5}} CH_ {3} COOH + C_ {2} H_ {5} OH} \end{array}
$

机理 2：

（图：）

请推测⼄酸⼄酯在超临界温度时的⽔解反应按上述哪种机理进⾏。为什么？

相关热⼒学数据如下：

表1 有关物质的标准热⼒学数据(25 °C)
表2 ⽔和⼄酸在25 MPa及不同温度下的解离常数 ${ \bf K } _ { \bf w }$ 和 $\lvert K _ { \mathrm { a } }$

\textit{（原卷此处含表格/仪器试剂表，详见 PDF 或 MD 源文件。）}

\textit{（原卷此处含表格/仪器试剂表，详见 PDF 或 MD 源文件。）}

\subsection*{5-4}
⼄酸⼄酯(EA)⽔解反应速率可表⽰为：

$
r = k c _ {\mathrm{EA}} \big (c _ {0, \mathrm{EA}} - c _ {\mathrm{EA}} \big) ^ {1 / 2}
$

其中�为速率常数， $c _ { 0 , \mathrm { E A } }$ 为⼄酸⼄酯的初始浓度。请通过推导说明⽔解反应按哪种机理进⾏。

\subsection*{5-5}
实验表明：近（超）临界⽔中酯类⽔解反应的表观活化能可降到常规条件下的 $1 / 2$ ，⽔解反应速率⼤幅度提⾼。请通过机理1分析原因。
